\begin{definition}[3.2.2]
\label{IV.2.3.2.2-fr}
Sean $X$ un preesquema localmente noetheriano, $\mathcal F$ un
$\mathcal O_X$-módulo coherente. Se dice que $\mathcal F$ es reducido
si satisface las dos condiciones siguientes~: 1\textsuperscript{o}
$\mathcal F$ carece de ciclo primo asociado inmerso~; 2\textsuperscript{o}
para todo punto maximal $x$ de $\operatorname{Supp}(\mathcal F)$, se tiene
$\operatorname{long}(\mathcal F_x)=1$.
\end{definition}

La condición 1\textsuperscript{o} significa que los ciclos primos asociados a
$\mathcal F$ son las componentes irreducibles de
$\operatorname{Supp}(\mathcal F)$ (3.1.4), y la condición
2\textsuperscript{o} significa que para todo punto genérico $x$ de una de tales
componentes se tiene $\operatorname{long}(\mathcal F_x)=1$.

Para un esquema afín $X$, esta definición da en particular la noción de
\emph{módulo reducido} sobre un anillo noetheriano $A$~; se dice que un
$A$-módulo de tipo finito $M$ es \emph{reducido} si no tiene ideales primos
asociados inmersos y si, para todo
$\mathfrak p\in\operatorname{Ass}(M)$,
$\operatorname{long}_{A_\mathfrak p}(M_\mathfrak p)=1$. Volviendo al caso de
un preesquema localmente noetheriano $X$ y de un $\mathcal O_X$-módulo
coherente $\mathcal F$, se dice que $\mathcal F$ es \emph{reducido en un
punto} $x\in X$ si $\mathcal F_x$ es un $\mathcal O_x$-\emph{módulo reducido}~; esto
significa asimismo que, sobre el esquema local
$\operatorname{Spec}(\mathcal O_x)$, $\widetilde{\mathcal F}_x$ es reducido~;
equivale, pues, a decir \emph{que $x$ no pertenece a ningún ciclo primo
asociado inmerso de $\mathcal F$ y que
$\operatorname{long}(\mathcal F_z)=1$ para todos los puntos maximales $z$ de
$\operatorname{Supp}(\mathcal F)$ tales que $x\in\overline{\{z\}}$}. Es claro
que si $\mathcal F$ es un $\mathcal O_X$-módulo coherente reducido, es
reducido en cada punto de $X$~; recíprocamente, si $\mathcal F$ es
\emph{reducido en el punto} $x$, existe una vecindad abierta $U$ de $x$ tal que
$\mathcal F|U$ sea un $\mathcal O_U$-módulo \emph{reducido}~: en efecto,
basta tomar $U$ que no interseque ningún ciclo primo asociado inmerso de
$\mathcal F$ (tal vecindad existe, pues esos ciclos forman un conjunto
localmente finito de partes
\oldpage[IV-2]{41}
cerradas de $X$). Decir que $\mathcal O_X$ es reducido en un punto $x$
equivale a decir que $X$ es \emph{reducido} en el punto $x$.
